\chapter{Conclusion} \label{ch:conclusion}

This thesis presented the first empirical comparison of four Bitcoin covenant vault
designs---CTV (BIP-119), CCV (BIP-443), OP\_VAULT (BIP-345), and CAT+CSFS
(BIP-347 + BIP-348)---plus a fifth vault built with the Simplicity language on
Elements. Through 15 experiments on regtest, Alloy formal verification, and
parameterized economic modeling, we established two results that constrain the
design space of covenant vaults.

\paragraph{Result 1: Griefing--safety impossibility.}
No covenant vault simultaneously achieves permissionless recovery and griefing
resistance (Proposition~\ref{thm:griefing_safety}). This is a structural design
constraint, not an implementation limitation. Vault designers must choose between
maximizing fund safety under key loss (CCV's keyless recovery) and minimizing
denial-of-service exposure (OP\_VAULT's authorized recovery).

\paragraph{Result 2: Structural design-space decomposition.}
Covenant vault designs are characterised by four nearly-orthogonal
dimensions: fee model ($f$), amount flexibility ($a$), recovery
gating ($g$), and recovery binding ($b$). Each dimension carries a
deterministic feature-to-vulnerability implication
(Lemmas~\ref{lem:anchor_pinning}--\ref{lem:unbound_theft}): the
anchor-output fee model opens TM1 (fee pinning); partial-withdrawal
flexibility opens TM4 (per-UTXO recovery-cost scaling); keyless recovery
opens TM2 (permissionless griefing, by
Proposition~\ref{thm:griefing_safety}); and unbound recovery opens
TM11 (cold-key theft for the full vault value).
Proposition~\ref{thm:fee_inversion} establishes that no proposed
Bitcoin Script extension occupies the dominant design-space corner
$(f_{\neq\text{anchor}}, a_{\text{atomic}}, g_{\text{key}},
b_{\text{bound}})$ at which $\mathcal{L} = 0$ simultaneously across
\{TM1, TM2, TM4, TM11\}. Every proposal differs from the corner in
at least one dimension, and by the corresponding lemma that
difference forces a positive defender loss on the threat that the
dimension controls. The decomposition also provides a framework for
localising future BIPs in the lattice: their vulnerability profile
follows from their feature tuple without further measurement.

\paragraph{Empirical corrections.}
We corrected prior hand-estimates of OP\_VAULT transaction sizes by 45--46\%
(trigger: $200 \rightarrow 292$\vB{}, recovery: $170 \rightarrow 246$\vB{}),
revealing a 36\% lifecycle cost premium over CCV. We measured CCV's watchtower
splitting capacity at approximately 6{,}172 splits per block---80\% more than
Harding's estimate of ${\sim}3{,}000$, which assumed OP\_VAULT-sized transactions.
These are the first regtest-validated vsize measurements for all four Bitcoin
covenant vault lifecycles.

\paragraph{Batched recovery and ordering flip.}
BIP-345's ``Batching'' section (authorised recovery, lines~613--621)
and BIP-443's default aggregation mode permit the defender to recover
$N$ Unvaulting UTXOs in a single transaction. Measured on regtest, the
amortised per-input cost is $68$\vB{} for CCV and $92$\vB{} for
OP\_VAULT. The dust-threshold fee rate rises by $1.8\times$ and
$2.7\times$ respectively, and the relative vulnerability of the two
designs under TM4 reverses: CCV is safer single-input, OP\_VAULT is
safer batched (Section~\ref{ss:batched_defender}).

\paragraph{Original implementations.}
We built a CAT+CSFS vault using explicit \texttt{OP\_CHECKSIGFROMSTACK}
(BIP-348) dual verification---distinct from prior OP\_CAT-only vaults
that use the Schnorr G-trick without CSFS~\cite{Poe21, Rij24}---and a
Simplicity vault on Elements using jet-based covenant enforcement.
Both are open-source. The CAT+CSFS implementation revealed the cold
key recovery vulnerability (TM11): bare \texttt{OP\_CHECKSIG} recovery
grants immediate, unrestricted fund theft upon cold key compromise.
The Simplicity vault serves as a cross-substrate reference point,
showing output-constrained recovery on all paths under the federated
sidechain trust model.

\paragraph{Formal verification.}
Alloy bounded model checking verified fund conservation, recovery
liveness, and no-extraction properties across the four Bitcoin Script
extension designs plus the Simplicity reference. The one expected
assertion failure---cold key extraction in CAT+CSFS---structurally
confirms TM11. The three-layer evidence methodology (Alloy for
structural possibility, regtest for concrete cost, fee model for
economic rationality) provides complementary assurance that no single
analysis technique could achieve alone.

\section{Future Work}

Several directions extend this research:

\begin{itemize}
\item \textbf{TRUC/v3 impact.} If adopted, TRUC transactions would eliminate CTV's
  descendant-chain fee pinning, removing its primary vulnerability and requiring
  re-evaluation of the TM1 branch of
  Proposition~\ref{thm:fee_inversion}'s counterexample enumeration.

\item \textbf{Cluster mempool.} The Bitcoin Core cluster mempool redesign may change
  descendant limit semantics, affecting the fee pinning attack construction.

\item \textbf{Temporal verification.} TLA+ or UPPAAL models could verify recovery
  race timing properties that Alloy's state-machine abstraction cannot capture.

\item \textbf{CTV+CSFS activation.} As CTV and CSFS move toward activation signaling,
  the measurements and cost models in this work can inform deployment decisions.
  The fee-conditional guidance (Section~\ref{se:deployment}) is directly applicable
  to vault designers choosing between CTV-only and CAT+CSFS vaults.

\item \textbf{Multi-vault composition.} Production deployments managing multiple
  concurrent vaults introduce fee wallet contention, cross-vault batching
  opportunities, and UTXO management complexity not captured by single-vault analysis.
\end{itemize}

The framework, all measurements, formal models, and both original vault
implementations are available at
\url{https://github.com/PraneethGunas/covenant-vault-research}.
